\section{Affine packets, parity necklaces, and exact residue automata}
\label{sec:affine-packets-residue-automata}

This section reconstructs an exponent-packet programme that arose in the
modern Chatnotes corpus.  The fixed-return formula itself is classical:
B\"ohm--Sontacchi and Crandall give parity-coordinate forms,
Lagarias proves the rational parity-word correspondence, and Urata writes the
odd-step exponent formula and asks when its rational periodic points are
integral
\cite{BohmSontacchi1978,Crandall1978,Lagarias1990,Urata2003}.
The comparison below identifies the rational points, transports cyclic
rotation together with its stabilizers, and relates the full, shortened,
and odd-return clocks.  It also resolves the packet label $(m,A)$ into an
exact fixed-content primitive-necklace count and a packetwise sign split.
The same exponent coordinates give a finite-state recursion for numerator
residues, retaining their multiplicities.

\subsection{Odd-step coordinates and the classical parity numerator}

Let $\mathbb Z_{(2)}=\{u/v\in\mathbb Q:u,v\in\mathbb Z,
v\text{ odd}\}$.  The parity of $u/v$ is $u\bmod2$; it is independent of
the odd-denominator representation.  Both maps
\[
 T(x)=\begin{cases}(3x+1)/2,&x\text{ odd},\\x/2,&x\text{ even},\end{cases}
 \qquad
 F(x)=\begin{cases}3x+1,&x\text{ odd},\\x/2,&x\text{ even}\end{cases}
\]
take $\mathbb Z_{(2)}$ to itself.  We call $T$ the shortened map and $F$
the full map.  Lagarias uses $T$ in his rational-cycle theorem
\cite[(1.1), Theorem~2.1]{Lagarias1990}.  Let
\[
 \mathcal U_2=\left\{\frac uv\in\mathbb Q:
     u,v\in\mathbb Z\text{ odd},\ v\ne0\right\}.
\]
For $x\in\mathcal U_2\setminus\{-1/3\}$, the numerator of $3x+1$
in odd-denominator form is even.  Thus
\[
 \Phi(x)=\frac{3x+1}{2^{\nu_2(3x+1)}}
 \quad\text{and}\quad
 a(x)=\nu_2(3x+1)\in\mathbb N_{>0}
\]
are well defined and $\Phi(x)\in\mathcal U_2$.

For $m\ge1$ and a packet $p=(a_1,\ldots,a_m)\in\mathbb N_{>0}^{m}$ put
\[
 A=A(p)=\sum_{i=1}^{m}a_i,
 \qquad S_0=0,\qquad S_j=\sum_{i=1}^{j}a_i,
\]
and define
\begin{equation}
 C(p)=\sum_{j=0}^{m-1}3^{m-1-j}2^{S_j},
 \qquad D_{m,A}=2^A-3^m,
 \qquad x_p=\frac{C(p)}{D_{m,A}}.
 \label{eq:packet-C-D-x}
\end{equation}
The equality $2^A=3^m$ is impossible, so $D_{m,A}\ne0$.  Moreover
$C(p)$ and $D_{m,A}$ are odd.  For tail recurrences only, set
$A(\varnothing)=C(\varnothing)=0$; no point $x_{\varnothing}$ is defined.

Write
\[
 \Pi_{m,A}=\{p\in\mathbb N_{>0}^{m}:A(p)=A\}.
\]
Here $A$ is a nonnegative integer, and $\Pi_{m,A}=\varnothing$ when
$A<m$.  For $A\ge m$ and $p\in\Pi_{m,A}$ define the shortened parity word
\begin{equation}
 E(p)=1,0^{a_1-1}1,0^{a_2-1}\cdots1,0^{a_m-1}
 \in\{0,1\}^{A}.
 \label{eq:packet-parity-encoding}
\end{equation}
If $v=(v_0,\ldots,v_{A-1})$ is a binary word, set
\begin{equation}
 N(v)=\sum_{j=0}^{A-1}v_j\,2^j
       3^{v_{j+1}+\cdots+v_{A-1}}.
 \label{eq:lagarias-numerator}
\end{equation}
This is the numerator in Lagarias's rational-cycle formula
\cite[Theorem~2.1]{Lagarias1990}.

\begin{theorem}[Packet--parity coordinate isomorphism]
\label{thm:packet-parity-coordinate-isomorphism}
For fixed $m\ge1$ and $A\ge m$, the map $E$ is a bijection
\[
 E:\Pi_{m,A}\longrightarrow
 \mathcal B_{m,A}^{\bullet}:=
 \{v\in\{0,1\}^{A}:v_0=1,
              \textstyle\sum_{j=0}^{A-1}v_j=m\}.
\]
It satisfies
\[
 N(E(p))=C(p).
\]
Consequently the rational point attached by Lagarias's formula to $E(p)$ is
exactly $x_p$.
\end{theorem}

\begin{proof}
The $m$ displayed blocks in \eqref{eq:packet-parity-encoding} have lengths
$a_1,\ldots,a_m$, hence total length $A$, and their initial symbols are
exactly the $m$ ones.  Conversely, starting at the distinguished one $v_0$,
record the positive cyclic distances to the successive ones and finally to
the boundary at $A$.  These distances form the unique inverse packet in
$\Pi_{m,A}$.

The ones of $E(p)$ occur at positions $S_0,\ldots,S_{m-1}$.  The number of
ones strictly to the right of $S_j$ is $m-1-j$.  Substitution in
\eqref{eq:lagarias-numerator} therefore gives
\[
 N(E(p))=\sum_{j=0}^{m-1}2^{S_j}3^{m-1-j}=C(p).
\]
Lagarias's denominator is $2^A-3^m$, so the two rational points coincide.
\end{proof}

The preceding bijection retains a distinguished odd point.  Advancing that
point by one odd return shifts the shortened word by $a_1$ positions.

Let $\sigma(a_1,\ldots,a_m)=(a_2,\ldots,a_m,a_1)$ and let $\rho$ be
left rotation of a binary word by one position.  The packet groupoid has
objects $p\in\Pi_{m,A}$ and arrows $(p,j):p\to\sigma^jp$ indexed by
$j\in\mathbb Z/m\mathbb Z$.  Distinct $j$ remain distinct arrows even
when their targets coincide.  The binary groupoid is the full subgroupoid
of the $\mathbb Z/A\mathbb Z$ rotation action whose objects belong to
$\mathcal B_{m,A}^{\bullet}$.  Its arrows $(v,t):v\to\rho^tv$ are
exactly those with $v_t=1$.

\begin{theorem}[Exact cyclic groupoid morphism]
\label{thm:packet-parity-cyclic-groupoid}
The object map $E$ and arrow map
\begin{equation}
 (p,j)\longmapsto\bigl(E(p),S_j(p)\bigr),
 \qquad 0\le j<m,
 \label{eq:packet-parity-arrow-map}
\end{equation}
give an isomorphism from the cyclic action groupoid of packets to the
pointed-one rotation groupoid of binary words.  Explicitly,
\begin{equation}
 E(\sigma^jp)=\rho^{S_j(p)}E(p).
 \label{eq:packet-parity-rotation}
\end{equation}
The inverse reads the cyclic gaps between consecutive ones.  In particular,
$p$ is primitive under packet rotation if and only if $E(p)$ is primitive
under binary rotation, and the induced map on necklace classes is bijective.
\end{theorem}

\begin{proof}
Moving the distinguished one past the first $j$ zero-blocks rotates the
binary word by their total length $S_j(p)$, which proves
\eqref{eq:packet-parity-rotation}.  For composable arrows,
\[
 S_{j+k\bmod m}(p)
 \equiv S_j(p)+S_k(\sigma^jp)\pmod A,
\]
so \eqref{eq:packet-parity-arrow-map} preserves composition.  The positions
of the ones are exactly $0=S_0<S_1<\cdots<S_{m-1}<A$.  Thus each target
arrow $(E(p),t)$ has the unique preimage $(p,j)$ with $t=S_j$.
The object inverse reads gaps, and this arrow inverse proves bijectivity on
every source fibre.  The identity maps to the identity.  For $j>0$,
\[
 S_{m-j}(\sigma^jp)=A-S_j(p),
\]
so inverses are preserved as well; $j=0$ is immediate.

Write $p=u^r$ with $u$ primitive, of length $d$ and sum $B$.
Then $m=rd$, $A=rB$, and the stabilizers are
\[
 \operatorname{Stab}(p)=\{kd\bmod m:0\le k<r\},\qquad
 \operatorname{Stab}(E(p))=\{kB\bmod A:0\le k<r\}.
\]
Indeed, a stabilizing binary rotation must send the initial one to a one,
hence is $S_j$ for some $j$.  The object bijection then forces
$\sigma^jp=p$, so $d\mid j$ and $S_{kd}=kB$.  Conversely all these
rotations stabilize $E(p)=E(u)^r$.  The arrow map sends $kd$ to $kB$.
In particular both stabilizers are trivial exactly when $r=1$; this proves
the primitive assertion and the bijection on cyclic classes.
\end{proof}

This is the precise interface between the exponent necklaces used by Urata
and the parity necklaces counted by Lagarias and later discussed through
Lyndon words by Laarhoven--de Weger
\cite{Lagarias1990,LaarhovenDeWeger2012,Urata2003}.

\subsection{Cyclic transport and the exact integrality criterion}

The local antecedent's Theorem~10.1 already proves affine cyclic transport
by uniqueness of the shifted fixed point.  The calculation below makes the
numerator identity, exact valuation, and two directions of divisibility
explicit in the same coordinates.

Put $p'=(a_2,\ldots,a_m)$.  Directly from
\eqref{eq:packet-C-D-x},
\[
 C(p)=3^{m-1}+2^{a_1}C(p'),
 \qquad
 C(\sigma p)=3C(p')+2^{A-a_1}.
\]

\begin{theorem}[Cyclic affine transport]
\label{thm:cyclic-affine-transport}
For every positive packet $p$,
\begin{equation}
 3C(p)+D_{m,A}=2^{a_1}C(\sigma p),
 \qquad
 3x_p+1=2^{a_1}x_{\sigma p}.
 \label{eq:cyclic-affine-transport}
\end{equation}
Hence
\[
 \nu_2(3x_p+1)=a_1,
 \qquad \Phi(x_p)=x_{\sigma p}.
\]
For fixed $m,A$, divisibility $D_{m,A}\mid C(p)$ is invariant under every
cyclic rotation of $p$.
\end{theorem}

\begin{proof}
Substitution gives
\[
 \begin{aligned}
 3C(p)+D_{m,A}
 &=3\bigl(3^{m-1}+2^{a_1}C(p')\bigr)+2^A-3^m\\
 &=2^{a_1}\bigl(3C(p')+2^{A-a_1}\bigr)
 =2^{a_1}C(\sigma p).
 \end{aligned}
\]
Division by the nonzero $D_{m,A}$ proves the rational identity.  Both
$C(\sigma p)$ and $D_{m,A}$ are odd, so $\nu_2(x_{\sigma p})=0$ and the
valuation is exactly $a_1$.  This also proves $3x_p+1\ne0$, since the
numerator $C(\sigma p)$ is positive.  Finally $\gcd(D_{m,A},6)=1$:
$D_{m,A}$ is odd and is congruent to $2^A\ne0$ modulo $3$.
In the first equality, invertibility of $2^{a_1}$ modulo $D_{m,A}$ gives
$D_{m,A}\mid C(p)\Rightarrow D_{m,A}\mid C(\sigma p)$; invertibility
of $3$ gives the converse.  Repeating this argument handles every rotation.
\end{proof}

\begin{corollary}[Typed cycle criterion]
\label{cor:typed-positive-cycle-criterion}
Let $p\in\Pi_{m,A}$.  Then $x_p$ is an integer if and only if
$D_{m,A}\mid C(p)$; whenever it is an integer it is odd.  It is positive if
and only if $D_{m,A}>0$.  If, in addition, $p$ is primitive, then
$x_p\in\mathbb Z_{>0}$ is equivalent to a positive odd integral orbit of
exact odd-step period $m$ with exponent necklace $[p]$.
\end{corollary}

\begin{proof}
The first assertion is the reduced rational integrality test before any
cancellation: $D_{m,A}\mid C(p)$.  An odd integer divided from odd $C(p)$ by
odd $D_{m,A}$ is odd.  Since $C(p)>0$, its sign is the sign of $D_{m,A}$.
Theorem~\ref{thm:cyclic-affine-transport} supplies every actual odd-step
arrow.  A shorter orbit period would make the exponent itinerary repeat,
contradicting primitivity.  Conversely, for $f_a(x)=(3x+1)/2^a$, induction
on the number of branches gives
\[
 (f_{a_m}\circ\cdots\circ f_{a_1})(x)
       =\frac{3^m x+C(p)}{2^A}.
\]
At an actual return this forces $D_{m,A}x=C(p)$, hence $x=x_p$.
If the itinerary were a proper power, its primitive root would itself
supply this periodic point by the same affine identity, contradicting the
assumed exact period.
\end{proof}

\begin{theorem}[Exact comparison of three clocks]
\label{thm:packet-three-clocks}
Let $p\in\Pi_{m,A}$ and $x_j=x_{\sigma^jp}$.  For $0\le j<m$,
\[
 \Phi^j(x_p)=T^{S_j}(x_p)=F^{S_j+j}(x_p)=x_j.
\]
Within its $(j+1)$st block the intermediate values are exactly
\begin{align*}
 T^s(x_j)&=2^{a_{j+1}-s}x_{j+1}
       &&(1\le s\le a_{j+1}),\\
 F^s(x_j)&=2^{a_{j+1}+1-s}x_{j+1}
       &&(1\le s\le a_{j+1}+1).
\end{align*}
The full parity word is
\[
 E_F(p)=1,0^{a_1}\cdots1,0^{a_m}\in\{0,1\}^{A+m}.
\]
If $p=u^r$ with $u$ primitive of length $d$ and sum $B$, the exact
periods of $x_p$ under $\Phi,T,F$ are respectively $d,B,B+d$.
\end{theorem}

\begin{proof}
The transport identity is $3x_j+1=2^{a_{j+1}}x_{j+1}$, with
$x_{j+1}$ odd and nonzero.  The first shortened step gives
$2^{a_{j+1}-1}x_{j+1}$; successive halving steps give the displayed
$T$ formula.  Its values are even before the last and odd at the last.
The full map inserts $2^{a_{j+1}}x_{j+1}$ as its first intermediate
value, proving the $F$ formula and its parity block.  Concatenation gives
all return times.  Thus $E_F$ inserts one zero after each one of $E$;
its inverse deletes that zero.  These are inverse bijections between
$\mathcal B_{m,A}^{\bullet}$ and the pointed binary words of length
$A+m$, weight $m$, having at least one zero between consecutive cyclic ones.

To establish exact periods, put $f_u=f_{u_d}\circ\cdots\circ f_{u_1}$.
It has positive slope $\alpha=3^d/2^B\ne1$, and its $r$th iterate is
$f_p$.  Since
\[
 f_u^r(x)-x=(1+\alpha+\cdots+\alpha^{r-1})(f_u(x)-x),
\]
the two maps have the same unique fixed point: $x_p=x_u$.
The least odd-return period is $d$, since a shorter period would repeat
the exponent itinerary of $u$.  Any return to the odd starting point under
$T$ or $F$ must be one of the odd returns just enumerated.  Their first
return to $x_u$ occurs after the $d$ blocks of $u$, at times $B$ and $B+d$.
This proves the three exact periods, including imprimitive packets.
\end{proof}

\subsection{Fixed-content primitive necklaces and the exact sign split}

For $m\geq1$ and $A\geq m$, let $W(m,A)$ be the number of words in
$\Pi_{m,A}$ whose least cyclic period is $m$.  Let
$N_{\mathrm{orb}}^{+}(m,A)$ and $N_{\mathrm{orb}}^{-}(m,A)$ denote the
numbers of positive and negative rational $\Phi$-orbits of exact odd-return
period $m$ represented by this packet, modulo cyclic phase.  Finally, let
$N_{\mathrm{pt}}^{\pm}(m,A)$ count the individual rational points in those
orbits.  These are three different counts: primitive pointed words, cyclic
orbits, and orbit points.

\begin{theorem}[Signed fixed-content necklace count]
\label{thm:signed-fixed-content-necklace-count}
Put $g=\gcd(m,A)$.  Then
\begin{equation}
 W(m,A)=\sum_{d\mid g}\mu(d)
            \binom{A/d-1}{m/d-1}.
 \label{eq:primitive-packet-word-count}
\end{equation}
The unsigned orbit count has the two exactly equal forms
\begin{equation}
 N_{\mathrm{orb}}(m,A)
 =\frac1m\sum_{d\mid g}\mu(d)
            \binom{A/d-1}{m/d-1}
 =\frac1A\sum_{d\mid g}\mu(d)
            \binom{A/d}{m/d}.
 \label{eq:primitive-packet-orbit-count}
\end{equation}
Its sign decomposition is
\begin{align}
 N_{\mathrm{orb}}^{+}(m,A)
 &=\mathbf 1_{\{2^A>3^m\}}N_{\mathrm{orb}}(m,A),
 &
 N_{\mathrm{orb}}^{-}(m,A)
 &=\mathbf 1_{\{2^A<3^m\}}N_{\mathrm{orb}}(m,A),
 \label{eq:signed-packet-orbit-count}\\
 N_{\mathrm{pt}}^{+}(m,A)
 &=mN_{\mathrm{orb}}^{+}(m,A),
 &
 N_{\mathrm{pt}}^{-}(m,A)
 &=mN_{\mathrm{orb}}^{-}(m,A).
 \label{eq:signed-packet-point-count}
\end{align}
There is no zero-sign case.  If $\gcd(m,A)=1$, every word in
$\Pi_{m,A}$ is primitive and therefore
\begin{equation}
 N_{\mathrm{orb}}(m,A)=\frac1m\binom{A-1}{m-1},
 \qquad
 N_{\mathrm{pt}}(m,A)=\binom{A-1}{m-1}.
 \label{eq:coprime-packet-orbit-count}
\end{equation}
\end{theorem}

\begin{proof}
There are $\binom{A-1}{m-1}$ positive compositions of $A$ into $m$
parts.  Every such word has a unique primitive root.  If it is the $d$th
power of a primitive word, then $d\mid m$ and $d\mid A$, and the root has
length $m/d$ and exponent sum $A/d$.  Conversely, repeating any primitive
word with those parameters $d$ times recovers a word in $\Pi_{m,A}$.
Consequently
\[
 \binom{A-1}{m-1}
 =\sum_{d\mid g}W(m/d,A/d).
\]
M\"obius inversion on the common divisors of $(m,A)$ gives
\eqref{eq:primitive-packet-word-count}; no Burnside coefficient remains
hidden in this inversion.

The cyclic action on a primitive word is free: a nonidentity stabilizing
rotation would exhibit the word as a proper power.  Each primitive necklace
therefore contains exactly $m$ pointed words.  Theorems
\ref{thm:packet-parity-cyclic-groupoid} and \ref{thm:packet-three-clocks}
identify one such necklace with one rational $\Phi$-orbit of exact period
$m$.  Conversely, the orbit determines its successive exact valuations, so
two necklaces cannot determine the same orbit.  This proves the first form
of \eqref{eq:primitive-packet-orbit-count}.  For each $d\mid g$,
\[
 \frac1m\binom{A/d-1}{m/d-1}
 =\frac1A\binom{A/d}{m/d},
\]
which proves the fixed-Hamming-weight binary-necklace form.

For every rotation $\sigma^jp$, equation \eqref{eq:packet-C-D-x} gives
\[
 x_{\sigma^jp}=\frac{C(\sigma^jp)}{2^A-3^m},
 \qquad C(\sigma^jp)>0.
\]
Thus all $m$ points of the orbit have the sign of $2^A-3^m$.  Equality
$2^A=3^m$ is impossible by unique factorization.  This proves
\eqref{eq:signed-packet-orbit-count}; exact period supplies the $m$ distinct
points in \eqref{eq:signed-packet-point-count}.  Finally, a proper power
would give a repetition number dividing both $m$ and $A$, so coprimality
forces every word to be primitive and yields
\eqref{eq:coprime-packet-orbit-count}.
\end{proof}

Stanley's multivariate Lyndon series counts primitive necklaces by the
number of occurrences of each colour
\cite[Exercise~7.89(a), equations~(7.149)--(7.150), and the solution on
printed p.~555]{Stanley2024}.  Extracting the binary-content coefficient
with content $(m,A-m)$ gives exactly the second expression in
\eqref{eq:primitive-packet-orbit-count}.  Lagarias gives the unrestricted
primitive binary-necklace and rational-cycle coordinates, while
Laarhoven--de Weger restate the unrestricted count in Lyndon-word language
\cite[pp.~38--39]{Lagarias1990}\cite[lines~303--315]{LaarhovenDeWeger2012}.
Thus the enumeration formula is classical.  The packet--parity groupoid
proves its exact transport to exponent compositions, and the denominator
$2^A-3^m$ supplies the Collatz-specific sign partition.  Summing over
$1\leq m\leq A$ recovers the unrestricted primitive binary-necklace count
for every $A\geq2$.  The sign refinement then uses the packet denominator
$2^A-3^m$ and does not alter the classical necklace enumeration.  At the
boundary, $m=1$ gives the negative fixed point for $A=1$ and one positive
fixed point for every $A\geq2$, whereas $A=m>1$ contains only the
imprimitive word $(1,\ldots,1)$ and contributes no exact-period-$m$ orbit.

\subsection{The exact minimum strip and its literature boundary}

Let $(x_0,\ldots,x_{m-1})$ be a positive odd integral cycle and write
$2^{a_i}x_{i+1}=3x_i+1$, with indices modulo $m$.  Multiplication around the
cycle gives the exact identity
\begin{equation}
 2^A=\prod_{i=0}^{m-1}\left(3+\frac1{x_i}\right).
 \label{eq:cycle-product-identity}
\end{equation}
This product identity and its minimum-point use are classical; compare
Crandall's cycle formula and minimum estimate
\cite[Section~7]{Crandall1978}.

If the cycle is nonconstant and $n=\min_i x_i$, then
\begin{equation}
 3^m<2^A<\left(3+\frac1n\right)^m.
 \label{eq:minimum-packet-strip}
\end{equation}
The left inequality is strict termwise.  The right inequality is strict
because at least one $x_i$ exceeds $n$.  A constant positive integral cycle
would satisfy $x=1/(2^a-3)$ and is therefore the fixed point $x=1$, $a=2$.
Moreover $1$ is fixed, while $3\mapsto5\mapsto1$ under $\Phi$; hence a
nontrivial positive cycle has $n\ge7$.  Thus
\begin{equation}
 3^m<2^A<\left(\frac{22}{7}\right)^m.
 \label{eq:seven-strip}
\end{equation}

No logarithmic approximation is required to enumerate the integral points in
this strip: test
\[
 3^m<2^A
 \quad\text{and}\quad
 7^m2^A<22^m
\]
as integer inequalities.  For $2\le m\le17$ the result is
\begin{center}
\begin{tabular}{c@{\qquad}l c@{\qquad}l}
\toprule
$m$ & admissible $A$ & $m$ & admissible $A$\\
\midrule
$2,3,4$ & none & $10$ & $16$\\
$5$ & $8$ & $11$ & $18$\\
$6,7$ & none & $12$ & none\\
$8$ & $13$ & $13$ & $21$\\
$9$ & none & $14$ & $23$\\
& & $15$ & $24$\\
& & $16$ & $26$\\
& & $17$ & $27,28$\\
\bottomrule
\end{tabular}
\end{center}

\begin{remark}[Historical comparison]
Crandall obtained a much larger cycle-length exclusion in 1978 from the
verification range available at that time \cite[Theorem~7.3]{Crandall1978}.
The small table here supplies explicit packet labels for the residue
calculations below.
\end{remark}

\subsection{Finite-state residue recursion}

For a positive odd modulus $q$, define
\[
 R_{m,A}(q)=\{C(p)\bmod q:p\in\Pi_{m,A}\}
 \subseteq\mathbb Z/q\mathbb Z.
\]
Writing $p=(a,p')$ gives
\begin{equation}
 C(a,p')=3^{m-1}+2^aC(p').
 \label{eq:first-letter-C-recursion}
\end{equation}

\begin{theorem}[Packet residue automaton]
\label{thm:packet-residue-automaton}
For $m\ge2$ and $A\ge m$,
\begin{equation}
 R_{m,A}(q)=
 \bigcup_{a=1}^{A-m+1}
 \left(3^{m-1}+2^aR_{m-1,A-a}(q)\right),
 \qquad R_{1,A}(q)=\{\bar1\}\quad(A\ge1).
 \label{eq:packet-residue-recursion}
\end{equation}
Let
\[
 v_{m,A}^{(q)}(r)=
 \#\{p\in\Pi_{m,A}:C(p)\equiv r\pmod q\}.
\]
For $A<m$ set $R_{m,A}(q)=\varnothing$ and $v_{m,A}^{(q)}=0$.
With column vectors, define
\[
 M_{m,a}^{(q)}e_r=e_{\,3^{m-1}+2^ar},\qquad
 (M_{m,a}^{(q)})^{-1}e_s=e_{\,2^{-a}(s-3^{m-1})},
\]
where all indices lie in $\mathbb Z/q\mathbb Z$.  Put
$\lambda=\operatorname{ord}_q(2)$ for $q>1$, and $\lambda=1$ for $q=1$.
Then the
formal vector series $V_m^{(q)}(z)=\sum_{A\ge m}v_{m,A}^{(q)}z^A$ obeys
\begin{equation}
 V_1^{(q)}(z)=\frac{z}{1-z}e_{\bar1},
 \qquad
 V_m^{(q)}(z)=
 \frac{\sum_{a=1}^{\lambda}z^aM_{m,a}^{(q)}}{1-z^{\lambda}}
 V_{m-1}^{(q)}(z).
 \label{eq:packet-residue-rational-series}
\end{equation}
Every coordinate of $V_m^{(q)}$ is therefore rational in $z$.
\end{theorem}

\begin{proof}
Every packet has a unique first part
$1\le a\le A-m+1$ and a unique tail in $\Pi_{m-1,A-a}$.
Equation~\eqref{eq:first-letter-C-recursion} proves the set recursion and,
with multiplicities retained, the vector recursion
\[
 v_{m,A}^{(q)}=\sum_{a=1}^{A-m+1}M_{m,a}^{(q)}v_{m-1,A-a}^{(q)}.
\]
Since $q$ is odd,
multiplication by $2^a$ is invertible modulo $q$, so each $M_{m,a}^{(q)}$ is
a permutation matrix.  The equality $2^{a+\lambda}\equiv2^a\pmod q$ makes
the matrix sequence $\lambda$-periodic.  Each $a\ge1$ has a unique form
$b+k\lambda$, $1\le b\le\lambda$, $k\ge0$.  In
$\mathbb Z^q[[z]]$, every coefficient of the following product is a finite
sum, and this decomposition gives
\[
 \sum_{a\ge1}z^aM_{m,a}^{(q)}
 =\frac{\sum_{b=1}^{\lambda}z^bM_{m,b}^{(q)}}{1-z^\lambda}.
\]
This proves \eqref{eq:packet-residue-rational-series}.  More explicitly,
putting $P_j(z)=\sum_{b=1}^{\lambda}z^bM_{j,b}^{(q)}$ gives
\[
 V_m^{(q)}(z)=
 \frac{P_m(z)P_{m-1}(z)\cdots P_2(z)\,z e_{\bar1}}
 {(1-z)(1-z^\lambda)^{m-1}}.
\]
The order of the matrix factors is as displayed.  For $m=1$ their product
is the identity.  When $q=1$ the formula reduces to
$V_m(z)=z^m/(1-z)^m$, the generating function for positive compositions.
\end{proof}

\begin{corollary}[Divisor exclusion]
\label{cor:packet-prime-factor-exclusion}
If $q$ is a positive divisor of $|D_{m,A}|$ and $0\notin R_{m,A}(q)$, then
$D_{m,A}\nmid C(p)$ for every $p\in\Pi_{m,A}$; hence the packet contains no
integral periodic point.
\end{corollary}

\begin{proof}
The divisibility $D_{m,A}\mid C(p)$ would imply $q\mid C(p)$, contradicting
$0\notin R_{m,A}(q)$.
\end{proof}

\subsection{Synchronized residues and the Chinese remainder map}

Factorwise information must retain the exponent word that produced each
residue.  The common-letter product of weighted automata is the standard
construction behind the Hadamard product of recognizable series
\cite[Definition~4.1.4, Proposition~4.3.10]{Wirth2022}.  Here its coordinate
map can be written explicitly.

\begin{theorem}[Synchronized CRT transport]\label{thm:packet-synchronized-crt}
Let $q_1,q_2$ be coprime positive odd integers and $q=q_1q_2$.
For each modulus $h$, let $L_h$ be the free abelian group with basis
$\{e_r:r\in\mathbb Z/h\mathbb Z\}$.  The map
\[
 J:L_q\longrightarrow L_{q_1}\otimes_{\mathbb Z}L_{q_2},\qquad
 Je_r=e_{r\bmod q_1}\otimes e_{r\bmod q_2}
\]
is an isomorphism preserving nonnegative coefficients.  Choose integers
$c_1,c_2$ with $q_2c_1\equiv1\pmod{q_1}$ and
$q_1c_2\equiv1\pmod{q_2}$.  Its inverse on basis vectors is
\[
 J^{-1}(e_u\otimes e_v)=e_{q_2c_1u+q_1c_2v\bmod q}.
\]
For every $m\ge2$ and $a\ge1$,
\begin{equation}\label{eq:crt-transition-intertwiner}
 JM_{m,a}^{(q)}=
 \bigl(M_{m,a}^{(q_1)}\otimes M_{m,a}^{(q_2)}\bigr)J.
\end{equation}
In particular, if
\[
 w_{m,A}(u,v)=\#\{p\in\Pi_{m,A}:
                  C(p)\equiv u\pmod{q_1},\ C(p)\equiv v\pmod{q_2}\},
\]
then $w_{m,A}=Jv_{m,A}^{(q)}$ and
\begin{equation}\label{eq:joint-residue-recursion}
 w_{m,A}=\sum_{a=1}^{A-m+1}
 \bigl(M_{m,a}^{(q_1)}\otimes M_{m,a}^{(q_2)}\bigr)w_{m-1,A-a}.
\end{equation}
The base vector is $w_{1,A}=e_{\bar1}\otimes e_{\bar1}$ for $A\ge1$,
and $w_{m,A}=0$ when $A<m$.
\end{theorem}

\begin{proof}
Reducing $q_2c_1u+q_1c_2v$ modulo either factor gives $u$ or $v$,
respectively.  If two integers have the same pair of residues, their
difference is divisible by both coprime factors and hence by $q$.
These facts prove both inverse identities, including independence from the
chosen representatives and Bezout coefficients.  Thus $J$ permutes the
basis, has zero kernel and singleton fibres, and preserves each coefficient.

Both sides of \eqref{eq:crt-transition-intertwiner} send $e_r$ to
\[
 e_{3^{m-1}+2^ar\bmod q_1}\otimes
 e_{3^{m-1}+2^ar\bmod q_2}.
\]
This proves the identity and, by composition, the corresponding identity
for every sequence of first-letter transitions.  Summing the basis image
over the same packets proves $w_{m,A}=Jv_{m,A}^{(q)}$.  Applying $J$ to
the scalar-modulus vector recursion proves
\eqref{eq:joint-residue-recursion}; no independent choice of letters in the
two factors occurs.  The base and empty ranges follow from their definitions.
\end{proof}

The factorwise inverse transition remains
$r\mapsto2^{-a}(r-3^{m-1})$.  Consequently this comparison loses no residue
information.  If $\lambda$ is the least common multiple of the two orders
of $2$ (with the singleton convention for modulus one), the joint generating
series $W_m(z)=\sum_{A\ge m}w_{m,A}z^A$ satisfies
\[
 W_m(z)=
 \frac{\sum_{a=1}^{\lambda}z^a
    (M_{m,a}^{(q_1)}\otimes M_{m,a}^{(q_2)})}{1-z^\lambda}
 W_{m-1}(z),\qquad
 W_1(z)=\frac z{1-z}e_{\bar1}\otimes e_{\bar1}.
\]
The proof is the same unique decomposition $a=b+k\lambda$ used above;
every coefficient is a finite sum.  Notice the single power $z^a$ and the
same letter $a$ in both factors.  A product of the separately summed
generating series would instead select two words and two budgets.
For any fixed $m,A$, the alphabet $1,\ldots,A$ and the length and budget
coordinates give finite layered automata.  Start at $(j,B,r)=(0,0,0)$,
read the reversed word $a_m,\ldots,a_1$, and use
\[
 (j,B,r)\longmapsto(j+1,B+a,3^j+2^ar\bmod q),
 \qquad j<m,\ B+a\le A.
\]
Accept at $j=m$, $B=A$.  The first transition sets the residue to $1$;
subsequent transitions prepend a letter.  The synchronized product shares
$j,B,a$ and retains the pair of residues.  It is the common-letter
construction cited above, restricted to the shared length and budget layers.

\begin{proposition}[Information discarded by the marginals]
\label{prop:packet-marginal-kernel}
On integer residue tables define
\[
 P(w)=\left(\left(\sum_v w_{uv}\right)_u,
                \left(\sum_u w_{uv}\right)_v\right).
\]
Its image consists of pairs of integer vectors with equal total sum.
Its kernel has the basis
\[
 E_{uv}-E_{u0}-E_{0v}+E_{00}
 \quad(u\ne0,\ v\ne0),
\]
where $E_{uv}=e_u\otimes e_v$, and has rank $(q_1-1)(q_2-1)$.
Each nonempty fibre is a translate of this kernel; for nonnegative count
tables it is the intersection of that translate with the nonnegative cone.
\end{proposition}

\begin{proof}
The two marginal totals both equal $\sum_{u,v}w_{uv}$.  Given marginals
$\alpha,\beta$ with common total $N$, the integer table
$\sum_u\alpha_uE_{u0}+\sum_v\beta_vE_{0v}-NE_{00}$ has those marginals.
Every displayed rectangle has zero row and column sums.  If $w$ has zero
marginals, subtract the sum of those rectangles with coefficients $w_{uv}$
for $u,v\ne0$.  The remaining table is supported on row zero and column
zero, and its zero marginals force every remaining entry to vanish.
Independence follows by inspecting each $(u,v)$ entry with $u,v\ne0$.
The fibre statements follow by subtracting two preimages.
\end{proof}

\begin{proposition}[Reduced denominators and integral $3x+d$ cycles]
\label{prop:packet-reduced-denominator-cycles}
Let $p\in\Pi_{m,A}$ with $D=2^A-3^m>0$, and set
$g=\gcd(C(p),D)$ and $d=D/g$.  Then $g$ is invariant under cyclic
rotation, and
\[
 n_j=\frac{C(\sigma^jp)}g=d\,x_{\sigma^jp}
\]
are positive odd integers satisfying $\gcd(n_j,d)=1$ and
\[
 3n_j+d=2^{a_{j+1}}n_{j+1},\qquad
 \nu_2(3n_j+d)=a_{j+1}.
\]
Multiplication by $d$ conjugates the rational shortened orbit to the
integral orbit of
\[
 T_d(n)=\begin{cases}(3n+d)/2,&n\text{ odd},\\n/2,&n\text{ even}.
 \end{cases}
\]
Its inverse on these orbits is division by $d$.  For primitive $p$, the
odd-return and shortened periods are exactly $m$ and $A$.
\end{proposition}

\begin{proof}
The identity $3C(p)+D=2^{a_1}C(\sigma p)$ and $\gcd(D,6)=1$
give $\gcd(C(p),D)=\gcd(C(\sigma p),D)$ by taking gcds with $D$.
Division by this common odd $g$ proves the displayed recurrence; its right
side has odd nonzero $n_{j+1}$, proving the valuation.  Reduction by the
full gcd gives $\gcd(n_j,d)=1$.  Since $d$ is odd, scaling by $d$
preserves parity, and checking the two branches gives
$T_d(dx)=dT(x)$ at every point of the rational orbit, including its even
intermediate points.  Division by $d$ proves both inverse identities and
preserves all return times.  The exact periods follow from
Theorem~\ref{thm:packet-three-clocks}.
\end{proof}

The source comparison here is Belaga--Mignotte's configuration-to-membership
map \cite[Sections~7--8]{BelagaMignotte2000}: their
$(k,\ell,a(P),B_{k,\ell},H,F,G)$ correspond respectively to
$(m,A,C(p),D,g,n_0,d)$.  Their aperiodic configuration is a primitive
exponent word; their primitive $3x+d$ cycle instead specifies coprimality
with $d$.  The two conditions serve different purposes in this correspondence.

\subsection{A complete joint-residue calculation at
\texorpdfstring{$(m,A)=(10,16)$}{(m,A)=(10,16)}}

Here
\[
 D_{10,16}=2^{16}-3^{10}=6487=13\cdot499,\qquad
 |\Pi_{10,16}|=\binom{15}{9}=5005.
\]
The CRT inverse is particularly explicit:
\[
 (u,v)\longmapsto3992u+2496v\pmod{6487},
 \qquad3992=499\cdot8,\quad2496=13\cdot192.
\]
The first coefficient is $(1,0)$ and the second $(0,1)$ in the two residue
coordinates.  Direct word enumeration and the independent joint recursion
give the following exact stratification:
\begin{center}
\begin{tabular}{r r r r}
\toprule
$\gcd(C(p),6487)$ & all words & primitive words & primitive cyclic classes\\
\midrule
$1$    & $4560$ & $4560$ & $456$\\
$13$   & $410$  & $410$  & $41$\\
$499$  & $35$   & $0$    & $0$\\
$6487$ & $0$    & $0$    & $0$\\
\midrule
Total  & $5005$ & $4970$ & $497$\\
\bottomrule
\end{tabular}
\end{center}
There are $3473$ occupied joint cells and $3014$ zero cells.  Both marginal
supports are full, but the marginal zero counts $410$ and $35$ have joint
zero count $w_{10,16}(0,0)=0$.  Thus a word realizing zero in one factor and
a word realizing zero in the other need not be the same word.

For explicit witnesses,
\[
 \begin{array}{c|r|r|r}
 p&C(p)&C(p)\bmod13&C(p)\bmod499\\\hline
 (1,1,1,1,1,1,2,2,2,4)&65065&0&195\\
 (1,1,1,1,4,1,1,1,1,4)&105289&2&0
 \end{array}
\]
The first word is primitive; the second is a square.  In fact every proper
repetition number divides both $10$ and $16$, so every imprimitive word is
$p=u^2$ for $u\in\Pi_{5,8}$.  There are exactly $35$ such words.  Expanding
the numerator in its first and second five-term blocks gives
\[
 C(u^2)=(3^5+2^8)C(u)=499C(u),\qquad
 D_{10,16}=(2^8-3^5)(2^8+3^5)=13\cdot499.
\]
The table shows that these are all the $499$-divisible words.  Therefore
$499$ alone rules out an integral periodic point for every primitive packet
at this label, although it does not exclude zero residues for all positive
words.  The $35$ squares form seven classes
of least odd period five, in addition to the $497$ primitive classes of
period ten.  The count of all cyclic classes is consequently $504$.

Under Proposition~\ref{prop:packet-reduced-denominator-cycles}, the first two
rows give $456$ coprime $3x+6487$ cycles and $41$ coprime $3x+499$ cycles,
each with ten odd returns and sixteen shortened steps.  The third row gives
seven coprime $3x+13$ cycles with five odd returns and eight shortened steps,
represented here twice.  The $456$ count is already reported by
Belaga--Mignotte \cite[Example~8.5(1)]{BelagaMignotte2000}; the $4560$
pointed words are its ten choices of odd base point per cycle.

These assertions are finite exact results: the supplied certificate checks
every one of the $5005$ words and every cell of the $13$-by-$499$ table.
Its joint recurrence uses the same first letter in both coordinates, and
its row and column sums agree with separately computed scalar recurrences.

\subsection{Exact finite packet certificates through period sixteen}

For $(m,A)=(5,8)$, $D_{5,8}=13$.  Since $\gcd(5,8)=1$, all $35$
packets are primitive and form seven cyclic classes.  The lexicographically
minimal representatives give:
\begin{center}
\begin{tabular}{c r r}
\toprule
$p$ & $C(p)$ & $C(p)\bmod13$\\
\midrule
$(1,1,1,1,4)$ & $211$ & $3$\\
$(1,1,1,2,3)$ & $227$ & $6$\\
$(1,1,1,3,2)$ & $259$ & $12$\\
$(1,1,2,1,3)$ & $251$ & $4$\\
$(1,1,2,2,2)$ & $283$ & $10$\\
$(1,1,3,1,2)$ & $331$ & $6$\\
$(1,2,1,2,2)$ & $319$ & $7$\\
\bottomrule
\end{tabular}
\end{center}
Theorem~\ref{thm:cyclic-affine-transport} makes divisibility invariant on
each class, so the packet has no integral point.

For $(m,A)=(8,13)$,
\[
 D_{8,13}=2^{13}-3^8=1631=7\cdot233.
\]
Here $|\Pi_{8,13}|=\binom{12}{7}=792$ and coprimality of $8$ and $13$
makes every packet primitive, hence there are $792/8=99$ cyclic classes.
The recursion \eqref{eq:packet-residue-recursion} gives the exact set
\begin{equation}
 R_{8,13}(233)=\mathbb Z/233\mathbb Z\setminus\{0,138\}.
 \label{eq:period-eight-residue-set}
\end{equation}
Among the $99$ minimal representatives, precisely $15$ have numerator
divisible by $7$, none has numerator divisible by $233$, and none is
integral.  Thus the factor $233$ excludes the entire packet.

The exact strip leaves the following candidate packets through period
sixteen.  Two independent streaming enumerations---separator combinations
with the displayed sum for $C$, and depth-first compositions with
$B_{j+1}=3B_j+2^{S_j}$---produce identical word streams and the following
counts.
\begin{center}
\begin{tabular}{r r r r r}
\toprule
$m$ & $A$ & $|\Pi_{m,A}|$ & $D_{m,A}$ & integral packets\\
\midrule
$5$  & $8$  & $35$      & $13$       & $0$\\
$8$  & $13$ & $792$     & $1631$     & $0$\\
$10$ & $16$ & $5005$    & $6487$     & $0$\\
$11$ & $18$ & $19448$   & $84997$    & $0$\\
$13$ & $21$ & $125970$  & $502829$   & $0$\\
$14$ & $23$ & $497420$  & $3605639$  & $0$\\
$15$ & $24$ & $817190$  & $2428309$  & $0$\\
$16$ & $26$ & $3268760$ & $24062143$ & $0$\\
\bottomrule
\end{tabular}
\end{center}
Each method processes exactly $4{,}734{,}620$ words.  In the packets
$(5,8)$, $(8,13)$, $(11,18)$, and $(13,21)$, the respective prime factors
\[
 13,\qquad233,\qquad7727,\qquad502829
\]
already exclude zero numerator residues.

\begin{theorem}[Certified finite consequence]
\label{thm:certified-period-sixteen-exclusion}
There is no nontrivial positive odd integral periodic orbit of exact odd-step
period $2\le m\le16$.
\end{theorem}

\begin{proof}
For $m=2,3,4,6,7,9,12$, the exact strip
\eqref{eq:seven-strip} contains no integer $A$.  The remaining values
$5,8,10,11,13,14,15,16$ have exactly the packets displayed in the table, and
each exhaustive divisibility count is zero.  The packet criterion in
Corollary~\ref{cor:typed-positive-cycle-criterion} completes the proof.
\end{proof}

The same strip first yields two packet labels at $m=17$, namely $(17,27)$
and $(17,28)$.  This is only the next finite locus selected by
\eqref{eq:seven-strip}; it is not asserted to be the frontier of the
published cycle literature.

\subsection{Reproducible calculations}

The portable script
\path{certificates/affine_packet_residue_checks.py} uses no floating point.
It checks the packet--parity numerator identity and all three clocks on
$12{,}089$ packets: $155{,}047$ shortened steps, $233{,}525$ full steps,
and $78{,}478$ odd returns.  These include $11{,}989$ primitive and $100$
imprimitive packets; the first return is checked in each clock.
On a separate $461$-packet range it checks the groupoid arrows, inverses,
composition, and stabilizers in both binary clocks.  Direct enumeration
checks $702$ residue multiplicity vectors, including $189$ empty ranges
$A<m$; $546$ coefficient vectors independently check the rational-series
recursion for the odd moduli $1,3,5,7,9,11,13$.  It computes
\eqref{eq:period-eight-residue-set} both by direct enumeration and by
\eqref{eq:packet-residue-recursion}.  Finally it runs the two independent
full scans above and pins every stream hash, minimum numerator, maximum
numerator, denominator, and count.  Optimized Python is rejected so that
assertions cannot be disabled.

The exhaustive scans stop at period sixteen; the two period-seventeen
labels are not included.  The general coordinate and generating-series
results are proved above independently of these finite checks.

The additional script \path{certificates/affine_crt_synchronization_checks.py}
independently checks the joint-residue calculation at $(10,16)$.  Two word
generators and three numerator formulas agree; it checks all $6487$ CRT
inverse pairs, $408681$ transition/inverse instances, the $35$ repeated-root
identities, and the full joint table including zero cells.  The certificate
rejects optimized Python and can write a new execution receipt without
overwriting an existing one.
